\documentclass[11pt,fontset=none]{ctexart}
% FandolSong is missing many traditional glyphs (e.g. 繐), so use Noto TC,
% which Overleaf has installed. Compile with XeLaTeX.
\setCJKmainfont{Noto Serif CJK TC}
\setCJKsansfont{Noto Sans CJK TC}
\setCJKmonofont{Noto Sans Mono CJK TC}

\usepackage[a4paper,landscape,margin=2cm]{geometry}
\usepackage{array,booktabs}
\usepackage{xeCJKfntef}   % \CJKunderline
\usepackage{xcolor}
\usepackage{fancyhdr}

% ---- footer ----
\pagestyle{fancy}
\fancyhf{}
\renewcommand{\headrulewidth}{0pt}
\fancyfoot[L]{\footnotesize 整理 @四分之一，参考衣若芬〈嚴謹的遊戲：王安石〈胡笳十八拍〉詩論析〉（2011），有改动。}
\fancyfoot[R]{\footnotesize\thepage}

\newcommand{\src}[1]{\CJKunderline{#1}}               % borrowed words
\newcommand{\unk}{\textcolor{gray}{未詳}}             % not found
\newcommand{\gap}{\textcolor{gray}{……}}               % missing half of a couplet

% ---- fixed column widths, so all pages line up ----
\newlength{\colA}\newlength{\colB}\newlength{\colC}\newlength{\colD}

% ---- one stanza per page, vertically centred ----
\newenvironment{stanzapage}[1]
  {\clearpage\vspace*{\stretch{1}}%
   \begin{center}%
   \begin{tabular}{@{}p{\colA}p{\colB}p{\colC}p{\colD}@{}}
   \toprule
   王安石《胡笳十八拍》 & 原作者 & 出處 & 原句\\
   \midrule
   \multicolumn{4}{@{}l}{\textbf{#1}}\\\addlinespace[2pt]}
  {\bottomrule
   \end{tabular}%
   \end{center}%
   \vspace*{\stretch{1.3}}}

\begin{document}

% widths are measured from the longest entry in each column
\settowidth{\colA}{如今豈無騕褭與驊騮，}
\settowidth{\colB}{南宮縣君錢氏}
\settowidth{\colC}{〈見王監兵馬使說近山有白黑二鷹〉之一}
\settowidth{\colD}{如今豈無騕褭與驊騮，時無王良伯樂死即休。}
\addtolength{\colC}{0.5em}
\addtolength{\colD}{0.5em}

\setlength{\tabcolsep}{1.6em}
\renewcommand{\arraystretch}{1.3}
\setlength{\defaultaddspace}{0.6em}

%% ---------------- 第一拍 ----------------
\begin{stanzapage}{第一拍}
中郎有女能傳業，
  & 韓愈 & 〈遊西林寺題蕭二兄郎中舊堂〉
  & \src{中郎有女能傳業}，伯道無兒可保家。\\
顏色如花命如葉。
  & 白居易 & 〈陵園妾‧憐幽閉也〉
  & 陵園妾，\src{顏色如花命如葉}。\\
\addlinespace
命如葉薄將奈何，
  & 白居易 & 〈陵園妾‧憐幽閉也〉
  & \src{命如葉薄將奈何}，一奉寢宮年月多。\\
一生抱恨常咨嗟。
  & 杜甫 & 〈負薪行〉
  & 更遭喪亂嫁不售，\src{一生抱恨}長\src{咨嗟}。\\
\addlinespace
良人持戟明光裏，
  & 張籍 & 〈節婦吟〉
  & 妾家高樓連苑起，\src{良人}執\src{戟明光裏}。\\
所慕靈妃媲簫史。
  & 韓愈 & 〈誰氏子〉
  & 或云欲學吹鳳笙，\src{所慕靈妃媲}蕭\src{史}。\\
\addlinespace
空房寂寞施繐帷，
  & 杜甫 & 〈過故斛斯校書莊二首〉之一
  & \src{空}餘\src{繐帷}在，淅淅野風秋。\\
棄我不待白頭時。
  & 張籍 & 〈白頭吟〉
  & 春天百草秋始衰，\src{棄我不待白頭時}。\\
\end{stanzapage}

%% ---------------- 第二拍 ----------------
\begin{stanzapage}{第二拍}
天不仁兮降亂離，
  & 蔡琰 & 〈胡笳十八拍〉第一拍
  & \src{天不仁兮降亂離}，地不仁兮使我逢此時。\\
嗟余去此其從誰。
  & 韓愈 & 〈祭田橫墓文〉
  & 死者不復生，\src{嗟余去此其從誰}。\\
\addlinespace
自胡之反持干戈，
  & 杜甫 & 〈寄柏學士林居〉
  & \src{自胡之反持干戈}，天下學士亦奔波。\\
翠蕤雲旓相蕩摩。
  & 杜甫 & 〈魏將軍歌〉
  & 欃槍熒惑不敢動，\src{翠蕤雲旓相蕩摩}。\\
\addlinespace
流星白羽腰間插，
  & 李白 & 〈塞上曲〉
  & \src{流星白羽腰間插}，劍花秋蓮光出匣。\\
疊鼓遙翻瀚海波。
  & 王維 & 〈燕支行〉
  & \src{疊鼓遙翻瀚海波}，鳴笳亂動天山月。\\
\addlinespace
一門骨肉散百草，
  & 李白 & 〈萬憤詞投魏郎中〉
  & \src{一門骨肉散百草}，遇難不復相提攜。\\
安得無淚如黃河。
  & 李商隱 & 〈安平公詩（故贈尚書韓氏）〉
  & 如公之德世一二，豈\src{得無淚如黃河}。\\
\end{stanzapage}

%% ---------------- 第三拍 ----------------
\begin{stanzapage}{第三拍}
身執略兮入西關，
  & 蔡琰 & 〈悲憤詩〉
  & \src{身執略兮入西關}，歷險阻兮之羌蠻。\\
關山阻修兮行路難。
  & 蔡琰 & 〈胡笳十八拍〉第十七拍
  & 十七拍兮心鼻酸，\src{關山阻修兮行路難}。\\
\addlinespace
水頭宿兮草頭坐，
  & 劉商 & 〈胡笳十八拍〉第五拍
  & \src{水頭宿兮草頭坐}，風吹漢地衣裳破。\\
在野只教心膽破。
  % full title: 〈見王監兵馬使說近山有白黑二鷹羅者久取竟未能得王以為毛骨有異他鷹恐臘後春生騫飛避暖勁翮思秋之甚眇不可見請余賦詩二首〉之一
  & 杜甫 & 〈見王監兵馬使說近山有白黑二鷹〉之一
  & \src{在野只教心}力\src{破}，於人何事網羅求。\\
\addlinespace
更鞴雕鞍教走馬，
  & 韋莊 & 〈秦婦吟〉
  & 旋抽金線學縫旗，才上\src{雕鞍教走馬}。\\
玉骨瘦來無一把。
  & 李商隱 & 〈偶成轉韻七十二句贈四同舍〉
  & 天官補吏府中趨，\src{玉骨瘦來無一把}。\\
\addlinespace
幾回拋鞚抱鞍橋，
  & 花蕊夫人 & 〈宮詞〉
  & 上得馬來才欲走，\src{幾回拋鞚抱鞍橋}。\\
往往驚墮馬蹄下。
  & 張籍 & 〈傷歌行〉
  & 郵夫防吏急喧驅，\src{往往驚墮馬蹄下}。\\
\end{stanzapage}

%% ---------------- 第四拍 ----------------
\begin{stanzapage}{第四拍}
漢家公主出和親，
  & 張籍 & 〈送和蕃公主〉
  & 塞上如今無戰塵，\src{漢家公主出和親}。\\
御廚絡繹送八珍。
  & 杜甫 & 〈麗人行〉
  & 黃門飛鞚不動塵，\src{御廚絡繹送八珍}。\\
\addlinespace
明妃初嫁與胡時，
  & 王安石 & 〈明妃曲二首〉之二
  & \src{明妃初嫁與胡}兒，氈車百兩皆胡姬。\\
一生衣服盡隨身。
  & 張籍 & 〈送和蕃公主〉
  & 九姓旗旛先引路，\src{一生衣服盡隨身}。\\
\addlinespace
眼長看地不稱意，
  & \unk & \unk & \unk\\
同是天涯淪落人。
  & 白居易 & 〈琵琶行〉
  & \src{同是天涯淪落人}，相逢何必曾相識。\\
\addlinespace
我今一食日還併，
  & 韓愈 & 〈寒食日出遊〉
  & 囊空甑倒誰救之，\src{我今一食日還併}。\\
短衣數挽不掩脛。
  & 杜甫 & 〈乾元中寓居同谷縣作歌七首〉之二
  & 黃獨無苗山雪盛，\src{短衣數挽不掩脛}。\\
\addlinespace
乃知貧賤別更苦，
  & 杜甫 & 〈醉歌行〉
  & \src{乃知貧賤別更苦}，吞聲躑躅涕淚零。\\
安得康強保天性。
  & 韓愈 & 〈寒食日出遊〉
  & 自然憂氣損天和，\src{安得康強保天性}。\\
\end{stanzapage}

%% ---------------- 第五拍 ----------------
\begin{stanzapage}{第五拍}
十三學得琵琶成，
  & 白居易 & 〈琵琶行〉
  & \src{十三學得琵琶成}，名屬教坊第一部。\\
繡幕重重卷畫屏。
  & \unk & \unk & \unk\\
\addlinespace
一見郎來雙眼明，
  & \unk & \unk & \unk\\
勸我酤酒花前傾。
  & 歐陽脩 & 〈啼鳥〉
  & 獨有花上提葫蘆，\src{勸我}沽\src{酒花前傾}。\\
\addlinespace
齊言此夕樂未央，
  & 張籍 & 〈楚宮行〉
  & 玉階羅幕微有霜，\src{齊言此夕樂未央}。\\
豈知此聲能斷腸。
  & 歐陽脩 & 〈明妃曲和王介甫作〉
  & 不識黃雲出塞路，\src{豈知此聲能斷腸}。\\
\addlinespace
如今正南看北斗，
  & 劉商 & 〈胡笳十八拍〉第六拍
  & 天翻地覆誰得知，\src{如今正南看北斗}。\\
言語傳情不如手。
  & 劉商 & 〈胡笳十八拍〉第六拍
  & 是非取與在指撝，\src{言語傳情不如手}。\\
\addlinespace
低眉信手續續彈，
  & 白居易 & 〈琵琶行〉
  & \src{低眉信手續續彈}，說盡心中無限事。\\
彈看飛鴻勸胡酒。
  & 王安石 & 〈明妃曲二首〉之二
  & 黃金捍撥春風手，\src{彈看飛鴻勸胡酒}。\\
\end{stanzapage}

%% ---------------- 第六拍 ----------------
\begin{stanzapage}{第六拍}
青天漫漫覆長路，
  & 張籍 & 〈送遠曲〉
  & \src{青天漫漫覆長路}，遠遊無家安得住。\\
一紙短書無寄處。
  & 晏殊 & 殘句（《庚溪詩話》作〈寄遠〉）
  & \src{一紙短書無寄處}，數行征雁入南雲。\\
\addlinespace
月下長吟久不歸，
  & 李白 & 〈金陵城西樓月下吟〉
  & \src{月下}沉\src{吟久不歸}，古來相接眼中稀。\\
當時還見雁南飛。
  & 陳陶 & 〈鄱陽秋夕〉
  & 今夜重開舊砧杵，\src{當時還見雁南飛}。\\
\addlinespace
彎弓射飛無遠近，
  & 劉商 & 〈胡笳十八拍〉第九拍
  & 髯胡少年能走馬，\src{彎弓射飛無遠近}。\\
青塚路邊南雁盡。
  & 李商隱 & 〈聞歌詩〉
  & \src{青塚路邊南雁盡}，細腰宮裏北人過。\\
\addlinespace
兩處音塵從此絕，
  & 蔡琰 & 〈胡笳十八拍〉第十拍
  & 故鄉隔兮\src{音塵絕}，哭無聲兮氣將咽。\\
唯向東西望明月。
  & 白居易 & 〈上陽白髮人〉
  & \src{唯向}深宮\src{望明月}，\src{東西}四五百回圓。\\
\end{stanzapage}

%% ---------------- 第七拍 ----------------
\begin{stanzapage}{第七拍}
明明漢月空相識，
  & 劉商 & 〈胡笳十八拍〉第四拍
  & 漫漫胡天叫不聞，\src{明明漢月}應\src{相識}。\\
道路只今多擁隔。
  & 杜甫 & 〈光祿坂行〉
  & 安得更似開元中，\src{道路}即\src{今多擁隔}。\\
\addlinespace
去住彼此無消息，
  & 杜甫 & 〈哀江頭〉
  & 清渭東流劍閣深，\src{去住彼此無消息}。\\
時獨看雲淚橫臆。
  & 杜甫 & 〈苦戰行〉
  & 別時孤雲今不飛，\src{時獨看雲淚橫臆}。\\
\addlinespace
豺狼喜怒難姑息，
  & 劉商 & 〈胡笳十八拍〉第二拍
  & 戎羯腥膻豈是人，\src{豺狼喜怒難姑息}。\\
自倚紅顏能騎射。
  & 杜甫 & 〈醉為馬所墜諸公攜酒相看〉
  & 向來皓首驚萬人，\src{自倚紅顏能騎射}。\\
\addlinespace
千言萬語無人會，
  & 鄭谷 & 〈燕詩〉
  & \src{千言萬語無人會}，又逐流鶯過短牆。\\
漫倚文章真末策。
  & \unk & \unk & \unk\\
\end{stanzapage}

%% ---------------- 第八拍 ----------------
\begin{stanzapage}{第八拍}
死生難有卻回身，
  & 李頻 & 〈太和公主還宮〉
  & 離亂應無初去貌，\src{死生難有卻回身}。\\
不忍重看舊寫真。
  & 王安石 & 〈明妃曲‧我本漢家子〉
  & 死生難有卻回身，\src{不忍}回\src{看舊寫真}。\\
\addlinespace
暮去朝來顏色改，
  & 白居易 & 〈琵琶行〉
  & 弟走從軍阿姨死，\src{暮去朝來顏色}故。\\
四時天氣總愁人。
  & 南宮縣君錢氏 & 詩（見《侯鯖錄》卷六）
  & 倘若有情相眷戀，\src{四時天氣總愁人}。\\
\addlinespace
東風漫漫吹桃李，
  & 王安石 & 〈金陵即事三首〉之二
  & \src{東風漫漫吹桃李}，非復當時仗外花。\\
盡日獨行春色裏。
  & 王安石 & 〈獨行〉
  & \src{盡日獨行春色裏}，醉吟誰肯伴衰翁。\\
\addlinespace
自經喪亂少睡眠，
  & 杜甫 & 〈茅屋為秋風所破歌〉
  & \src{自經喪亂少睡眠}，長夜沾濕何由徹。\\
鶯飛燕語長悄然。
  & 白居易 & 〈上陽白髮人〉
  & \src{鶯}歸\src{燕}去\src{長悄然}，春往秋來不記年。\\
\end{stanzapage}

%% ---------------- 第九拍 ----------------
\begin{stanzapage}{第九拍}
柳絮已將春去遠，
  % full title: 〈初至潁州西湖種瑞蓮黃楊寄淮南轉運呂度支發運許主客〉
  & 歐陽脩 & 〈初至潁州西湖種瑞蓮黃楊……〉
  & \src{柳絮已將春去遠}，海棠應恨我來遲。\\
攀條弄芳畏晼晚。
  & 王安石 & 〈移桃花示俞秀老〉
  & \src{攀條弄芳畏晼晚}，已見黍雪盤中毛。\\
\addlinespace
憂患眾兮歡樂鮮，
  & 潘岳 & 〈哀永逝文〉
  & 逝日長兮生年淺，\src{憂患眾兮歡樂鮮}。\\
一去可憐終不返。
  & 王安石 & 〈謝安墩二首〉之二
  & \src{一去可憐終不返}，暮年垂淚對桓伊。\\
\addlinespace
日夕思歸不得歸，
  & 劉商 & 〈胡笳十八拍〉第八拍
  & 旦\src{夕思歸不得歸}，愁心想似籠中鳥。\\
山川滿目淚沾衣。
  & 李嶠 & 〈汾陰行〉
  & \src{山川滿目淚沾衣}，富貴榮華能幾時。\\
\addlinespace
罼圭苑裏西風起，
  & 杜牧 & 〈故洛陽城有感〉
  & 篳\src{圭苑裏}秋\src{風}後，平樂館前斜日時。\\
歎息人間萬事非。
  & 杜甫 & 〈送韓十四江東省覲〉
  & 兵戈不見老萊衣，\src{歎息人間萬事非}。\\
\end{stanzapage}

%% ---------------- 第十拍 ----------------
\begin{stanzapage}{第十拍}
寒聲一夜傳刁斗，
  & 高適 & 〈燕歌行〉
  & 殺氣三時作陣雲，\src{寒聲一夜傳刁斗}。\\
雲雪埋山蒼兕吼。
  & 杜甫 & 〈復陰〉
  & 江濤簸岸黃沙走，\src{雲雪埋山蒼兕吼}。\\
\addlinespace
詩成吟詠轉淒涼，
  & 杜甫 & 〈至後〉
  & 愁極本憑詩遣興，\src{詩成吟詠轉淒涼}。\\
不如獨坐空搔首。
  & 高適 & 〈九月九日酬顏少府〉
  & 縱使登高只斷腸，\src{不如獨坐空搔首}。\\
\addlinespace
漫漫胡天叫不聞，
  & 劉商 & 〈胡笳十八拍〉第四拍
  & \src{漫漫胡天叫不聞}，明明漢月應相識。\\
胡人高鼻動成群。
  & 杜甫 & 〈黃河二首〉之一
  & 鐵馬長鳴不知數，\src{胡人高鼻動成群}。\\
\addlinespace
寒盡春生洛陽殿，
  % full title: 〈湖城東遇孟雲卿復歸劉顥宅宿宴飲散因為醉歌〉
  & 杜甫 & 〈湖城東遇孟雲卿復歸劉顥宅……〉
  & 天開地裂長安陌，\src{寒盡春生洛陽殿}。\\
回首何時復來見。
  & 司馬光 & 〈和王介甫明妃曲〉
  & 宮門銅環雙獸面，\src{回首何時復來見}。\\
\end{stanzapage}

%% ---------------- 第十一拍 ----------------
\begin{stanzapage}{第十一拍}
晚來幽獨恐傷神，
  & 杜甫 & 〈題鄭縣亭子〉
  & 更欲題詩滿青竹，\src{晚來幽獨恐傷神}。\\
唯見沙蓬水柳春。
  & 張籍 & 〈送和蕃公主〉
  & 氈城南望無回日，空\src{見沙蓬水柳春}。\\
\addlinespace
破除萬事無過酒，
  & 韓愈 & 〈贈鄭兵曹〉
  & 杯行到君莫停手，\src{破除萬事無過酒}。\\
虜酒千盃不醉人。
  & 高適 & 〈營州歌〉
  & \src{虜酒千}鍾\src{不醉人}，胡兒十歲能騎馬。\\
\addlinespace
含情欲說更無語，
  & 朱慶餘 & 〈宮詞〉
  & \src{含情欲說}宮中事，鸚鵡前頭不敢言。\\
一生長恨奈何許。
  & 韓愈 & 〈感春四首〉之一
  & 三杯取醉不復論，\src{一生長恨奈何許}。\\
\addlinespace
饑對酪肉兮不能餐，
  & 蔡琰 & 〈胡笳十八拍〉第六拍
  & 冰霜凜凜兮身苦寒，\src{饑對}肉酪\src{兮不能餐}。\\
強來前帳臨歌舞。
  & 儲光羲 & 〈明妃曲〉
  & \src{強來前帳}看\src{歌舞}，共待單于夜獵歸。\\
\end{stanzapage}

%% ---------------- 第十二拍 ----------------
\begin{stanzapage}{第十二拍}
歸來展轉到五更，
  & 李商隱 & 〈無題四首〉之四
  & \src{歸來展轉到五更}，梁間燕子聞長歎。\\
起看北斗天未明。
  & 張籍 & 〈秋夜長〉
  & 荒城為村無更聲，\src{起看北斗天未明}。\\
\addlinespace
秦人築城備胡處，
  & 李白 & 〈戰城南〉
  & \src{秦}家\src{築城備胡處}，漢家還有烽火燃。\\
擾擾唯有牛羊聲。
  & 張籍 & 〈將軍行〉
  & 胡兒殺盡陰磧暮，\src{擾擾唯有牛羊聲}。\\
\addlinespace
萬里飛蓬映天過，
  & 杜甫 & 〈復陰〉
  & \src{萬里飛蓬映天過}，孤城樹羽揚風直。\\
風吹漢地衣裳破。
  & 劉商 & 〈胡笳十八拍〉第五拍
  & 水頭宿兮草頭坐，\src{風吹漢地衣裳破}。\\
\addlinespace
欲往城南望城北，
  & 杜甫 & 〈哀江頭〉
  & 黃昏胡騎塵滿城，\src{欲往城南望城北}。\\
三步回頭五步坐。
  & 杜甫 & 〈憶昔行〉
  & 千崖無人萬壑靜，\src{三步回頭五步坐}。\\
\end{stanzapage}

%% ---------------- 第十三拍 ----------------
\begin{stanzapage}{第十三拍}
自斷此生休問天，
  & 杜甫 & 〈曲江三章章五句〉之三
  & \src{自斷此生休問天}，杜曲幸有桑麻田。\\
生得胡兒擬棄捐。
  & 劉商 & 〈胡笳十八拍〉第十拍
  & \src{生得胡兒}欲\src{棄捐}，及生母子情宛然。\\
\addlinespace
一始扶床一初坐，
  & 白居易 & 〈母別子〉
  & \src{一始扶}行\src{一初坐}，坐啼行哭牽人衣。\\
抱攜撫視皆可憐。
  % full title: 〈送師厚歸南陽會天大風遂宿高陽山寺明日同至姜店〉
  & 梅堯臣 & 〈送師厚歸南陽會天大風……〉
  & 像塑神母乳九子，\src{抱攜撫}玩\src{皆可憐}。\\
\addlinespace
寧知遠使問名姓，
  & 劉商 & 〈胡笳十八拍〉第十二拍
  & \src{寧知遠使問}姓名，漢語泠泠傳好音。\\
引袖拭淚悲且慶。
  % full title: 〈寒食日出遊夜歸張十一院長見示病中憶花九篇因此投贈〉
  & 韓愈 & 〈寒食日出遊〉
  & 豈料生還得一處，\src{引袖拭淚悲且慶}。\\
\addlinespace
悲莫悲兮生別離，
  & 屈原 & 〈九歌‧大司命〉
  & \src{悲莫悲兮生別離}，樂莫樂兮新相知。\\
悲在君家留二兒。
  & 白居易 & 〈母別子〉
  & 迎新棄舊未足悲，\src{悲在君家留}兩\src{兒}。\\
\end{stanzapage}

%% ---------------- 第十四拍 ----------------
\begin{stanzapage}{第十四拍}
鞠之育之不羞恥，
  & 蔡琰 & 〈胡笳十八拍〉第十一拍
  & \src{鞠之育之}兮\src{不羞恥}，愍之念之兮生長邊鄙。\\
恩情亦各言其子。
  & 劉商 & 〈胡笳十八拍〉第十四拍
  & 莫以胡兒可羞恥，\src{恩情亦各言其子}。\\
\addlinespace
天寒日暮山谷裏，
  & 杜甫 & 〈乾元中寓居同谷縣作歌七首〉之一
  & 歲拾橡栗隨狙公，\src{天寒日暮山谷裏}。\\
腸斷非關隴頭水。
  & 李白 & 〈猛虎行〉
  & \src{腸斷非關隴頭水}，淚下不為雍門琴。\\
\addlinespace
兒呼母兮啼失聲，
  & 蔡琰 & 〈悲憤詩〉
  & \src{兒呼母兮啼失聲}，我掩耳兮不忍聽。\\
依然離別難為情。
  & 韓愈 & 〈桃源圖〉
  & 人間有累不可住，\src{依然離別難為情}。\\
\addlinespace
灑血仰頭兮訴蒼蒼，
  & 蔡琰 & 〈胡笳十八拍〉第十六拍
  & 泣\src{血仰頭兮訴蒼蒼}，胡為生我兮獨罹此殃。\\
知我如此兮不如無生。
  & 蔡琰 & 〈胡笳十八拍〉第三拍
  & 越漢國兮入胡城，亡家失身\src{兮不如無生}。\\
\end{stanzapage}

%% ---------------- 第十五拍 ----------------
\begin{stanzapage}{第十五拍}
當時悔來歸又恨，
  & 劉商 & 〈胡笳十八拍〉第十五拍
  & 歎息襟懷無定分，\src{當時}怨\src{來歸又恨}。\\
洛陽宮殿焚燒盡。
  & 杜甫 & 〈憶昔二首〉之二
  & \src{洛陽宮殿}燒焚\src{盡}，宗廟新除狐兔穴。\\
\addlinespace
紛紛黎庶逐黃巾，
  % full title: 〈自蘇臺至望亭驛人家盡空春物增思悵然有作因寄從弟紓〉
  & 李嘉祐 & 〈自蘇臺至望亭驛人家盡空……〉
  & 南浦菰蔣覆白蘋，東吳\src{黎庶逐黃巾}。\\
心折此時無一寸。
  & 杜甫 & 〈冬至〉
  & \src{心折此時無一寸}，路迷何處是三秦。\\
\addlinespace
慟哭秋原何處村，
  & 杜甫 & 〈白帝〉
  & 哀哀寡婦誅求盡，\src{慟哭秋原何處村}。\\
千家今有百家存。
  & 杜甫 & 〈白帝〉
  & 戎馬不如歸馬逸，\src{千家今有百家存}。\\
\addlinespace
爭持酒食來相饋，
  & 韓愈 & 〈桃源圖〉
  & \src{爭持酒食來相饋}，禮數不同樽俎異。\\
舊事無人可共論。
  & 韓愈 & 〈過始興江口感懷〉
  & 目前百口還相逐，\src{舊事無人可共論}。\\
\end{stanzapage}

%% ---------------- 第十六拍 ----------------
\begin{stanzapage}{第十六拍}
此身飲罷無歸處，
  & 杜甫 & 〈樂遊園歌〉
  & \src{此身飲罷無歸處}，獨立蒼茫自詠詩。\\
心懷百憂復千慮。
  & 杜甫 & 〈人日寄杜二拾遺〉
  & 身在南蕃無所預，\src{心懷百憂復千慮}。\\
\addlinespace
天翻地覆誰得知，
  & 劉商 & 〈胡笳十八拍〉第六拍
  & \src{天翻地覆誰得知}，如今正南看北斗。\\
魏公垂淚嫁文姬。
  & 舒元輿 & 〈贈李翺〉
  & 誰是蔡邕琴酒客，\src{魏公}懷舊\src{嫁文姬}。\\
\addlinespace
天涯憔悴身，
  & 崔塗 & 〈蜀城春望〉
  & \src{天涯憔悴身}，一望一沾巾。\\
託命於新人。
  & 蔡琰 & 〈悲憤詩〉
  & \src{託命於新人}，竭心自勖厲。\\
\addlinespace
念我出腹子，
  & 蔡琰 & 〈悲憤詩〉
  & \src{念我出腹子}，匈臆為摧敗。\\
使我歎恨勞精神。
  & 杜甫 & 〈苦戰行〉
  & 干戈未定失壯士，\src{使我歎恨}傷\src{精}魂。\\
\addlinespace
新人新人聽我語，
  & 白居易 & 〈母別子〉
  & \src{新人新人聽我語}，洛陽無限紅樓女。\\
我所思兮在何所。
  & 韓愈 & 〈感春四首〉之一
  & \src{我所思兮在何所}，情多地遐兮遍處處。\\
\addlinespace
母子分離兮意難任，
  & 蔡琰 & 〈胡笳十八拍〉第十五拍
  & 日月無私兮曾不照臨，子母\src{分離兮意難任}。\\
死生不相知兮何處尋。
  & 張籍 & 〈各東西〉
  & 我今與子非一身，安得\src{死生不相}棄。\\
\end{stanzapage}

%% ---------------- 第十七拍 ----------------
\begin{stanzapage}{第十七拍}
燕山雪花大如席，
  & 李白 & 〈北風行〉
  & \src{燕山雪花大如席}，片片吹落軒轅臺。\\
與兒洗面作光澤。
  & 北齊崔氏 & 〈靧面辭〉
  & 取花紅，取雪白，\src{與兒洗面作光澤}。\\
\addlinespace
怳然天地半夜白，
  & 歐陽脩 & 〈晏太尉西園賀雪歌〉
  & \src{怳然天地半夜白}，群雞失曉不及鳴。\\
閨中只是空相憶。
  & 岑參 & 〈題苜蓿峰寄家人〉
  & \src{閨中只是空}思想，不見沙場愁殺人。\\
\addlinespace
點注桃花舒小紅，
  & 杜甫 & 〈江雨有懷鄭典設〉
  & 寵光蕙葉與多碧，\src{點注桃花舒小紅}。\\
與兒洗面作華容。
  & 北齊崔氏 & 〈靧面辭〉
  & 取雪白，取花紅，\src{與兒洗面作華容}。\\
\addlinespace
欲問平安無使來，
  & 杜甫 & 〈所思〉
  & 可憐懷抱向人盡，\src{欲問平安無使來}。\\
桃花依舊笑春風。
  & 崔護 & 〈題都城南莊〉
  & 人面不知何處去，\src{桃花依舊笑春風}。\\
\end{stanzapage}

%% ---------------- 第十八拍 ----------------
\begin{stanzapage}{第十八拍}
春風似舊花仍笑，
  & 范仲淹 & 〈靈巖寺〉
  & \src{春風似舊花}猶\src{笑}，往事多遺石不言。\\
人生豈得長年少。
  & 李商隱 & 〈無題〉
  & \src{人生豈得長}無謂，懷古思鄉共白頭。\\
\addlinespace
我與兒兮各一方，
  & 蔡琰 & 〈胡笳十八拍〉第十六拍
  & 十六拍兮思茫茫，\src{我與兒兮各一方}。\\
憔悴看成兩鬢霜。
  & \unk & \unk & \unk\\
\addlinespace
如今豈無騕褭與驊騮，
  & 杜甫 & 〈天育驃騎歌〉
  & \src{如今豈無騕褭與驊騮}，時無王良伯樂死即休。\\
安得送我置汝傍。
  & 杜甫 & 〈乾元中寓居同谷縣作歌七首〉之三
  & 東飛駕鵝後鶖鶬，\src{安得送我置汝傍}。\\
\addlinespace
胡塵暗天道路長，
  & 杜甫 & 〈乾元中寓居同谷縣作歌七首〉之二
  & 生別展轉不相見，\src{胡塵暗天道路長}。\\
遂令再往之計墮眇芒。
  % full title: 〈盧郎中雲夫寄示送盤谷子詩兩章歌以和之〉
  & 韓愈 & 〈盧郎中雲夫寄示送盤谷子詩……〉
  & 歸來辛苦欲誰為，坐\src{令再往之計墮眇}茫。\\
\addlinespace
胡笳本出自胡中，
  & 蔡琰 & 〈胡笳十八拍〉第十八拍
  & \src{胡笳本}自出\src{胡中}，緣琴翻出音律同。\\
此曲哀怨何時終。
  & 杜甫 & 〈歲晏行〉
  & 萬國城頭吹畫角，\src{此曲哀怨何時終}。\\
\addlinespace
笳一會兮琴一拍，
  & 蔡琰 & 〈胡笳十八拍〉第一拍
  & \src{笳一會兮琴一拍}，心憤怨兮無人知。\\
此心炯炯君應識。
  & 杜甫 & 〈偪仄行贈畢曜〉
  & 徒步翻愁官長怒，\src{此心炯炯君應識}。\\
\end{stanzapage}

\end{document}
